\documentclass[12pt]{article}

% ------------------------------------------------------------
% Pacotes básicos
% ------------------------------------------------------------
\usepackage[T1]{fontenc}
\usepackage[utf8]{inputenc}
%\usepackage[brazil]{babel}
\usepackage[english]{babel}

\usepackage{lipsum}         % texto exemplo
\usepackage{graphicx}       % imagens
\usepackage{float}          % controle de posição de figuras
\usepackage{booktabs}       % tabelas
\usepackage{amsmath,amssymb}% matemática
\usepackage{todonotes}
\usepackage{hyperref}       % links clicáveis
\usepackage{caption}        % legendas personalizadas
\hypersetup{
  colorlinks   = true,
  urlcolor     = blue,
  linkcolor    = blue,
  citecolor    = blue
}
\usepackage{tikz}            % diagramas
\usepackage{microtype}       % melhorias tipográficas
\usepackage{enumitem}
\usetikzlibrary{patterns.meta, arrows.meta, positioning, shapes.geometric}

% ------------------------------------------------------------
% Todos and notes
% ------------------------------------------------------------
%\setlength{\marginparwidth}{2cm}
%\usepackage[colorinlistoftodos,textsize=scriptsize]{todonotes}

% ------------------------------------------------------------
% Bibliografia e Citações
% ------------------------------------------------------------
\usepackage{csquotes} % citações
\usepackage[backend=biber, style=abnt]{biblatex}
\addbibresource{references.bib}

% ------------------------------------------------------------
% Configurações do documento
% ------------------------------------------------------------
% Page setup
\usepackage{geometry}
\geometry{a4paper, margin=2.5cm}

% Paragraph spacing
\usepackage{parskip}

% Configuração dos títulos dos apêndices
\usepackage{titlesec}
\renewcommand{\appendixname}{APÊNDICE}

% ------------------------------------------------------------
% Variáveis do Documento
% ------------------------------------------------------------

\newcommand{\instituicao}{UNIVERSIDADE FEDERAL DE SANTA CATARINA}
\newcommand{\departamento}{Programa de Pós-Graduação em Economia}
\newcommand{\curso}{Mestrado}
\newcommand{\disciplina}{Macroeconomia I}
\newcommand{\titulo}{Modelo de Solow Swan}
\newcommand{\subtitulo}{Conceitos Preliminares}
\newcommand{\autor}{}
%\newcommand{\orientador}{}
\newcommand{\local}{Florianópolis}
\newcommand{\dataentrega}{\today}
%\newcommand{\resumo}{
%\lipsum[1][1-8] % Resumo do trabalho
%}

% ------------------------------------------------------------
% Título simples
% ------------------------------------------------------------
\newcommand{\imprimirtitulo}{
\noindent

\begin{center}
  {\Large\bfseries\MakeUppercase{\instituicao}\par}
  \vspace{0.2em}
  {\large \departamento\par}
  \vspace{0.2em}
  {\normalsize\textbf{\disciplina}\ (\curso)\par}
  \vspace{1.4em}
  {\LARGE\bfseries \titulo\par}
  \vspace{0.25em}
  {\large\itshape \subtitulo\par}
  \vspace{1.5em}
\end{center}

\noindent
\begin{minipage}{0.48\textwidth}
  \raggedright\normalsize \autor
\end{minipage}\hfill
\begin{minipage}{0.48\textwidth}
  \raggedleft\normalsize \local, \dataentrega
\end{minipage}

\vspace{1.0em}
\hrule
}


% ------------------------------------------------------------
% Capa
% ------------------------------------------------------------

\begin{document}
\imprimirtitulo

% ----------------------------------------------------------------------------
% Sumário
% ----------------------------------------------------------------------------

\renewcommand{\contentsname}{Sumário}
\tableofcontents
\vspace{1.5em}
\hrule

% ----------------------------------------------------------------------------
% Corpo do documento
% ----------------------------------------------------------------------------
\newpage
\section{Conceitos preliminares}

\subsection{Funções côncavas}

\textbf{Concavidade} $\Rightarrow$ \textbf{rendimentos marginais decrescentes} $\Rightarrow f''(x) < 0$. 

Geometricamente, a reta entre dois pontos da curva fica abaixo da função.

\noindent\rule{\textwidth}{0.4pt}

Uma função $f : D \subseteq \mathbb{R} \to \mathbb{R}$ é \textbf{côncava} se, para quaisquer $x,y \in D$ e qualquer $\lambda \in [0,1]$,
$f(\lambda x + (1-\lambda)y) \geq \lambda f(x) + (1-\lambda)f(y)$.

Em palavras, o valor da função na média ponderada entre $x$ e $y$ é maior ou igual à média ponderada dos valores da função nesses pontos. Geometricamente, isso significa que o gráfico da função fica acima de qualquer corda ligando dois pontos do seu gráfico.

Se a desigualdade for estrita para $x \neq y$ e $\lambda \in (0,1)$, a função é \textbf{estritamente côncava}.

Se $f$ for diferenciável, a concavidade implica que \textbf{a primeira derivada $f'$} é decrescente. Isto é, se $x_1 < x_2$, então
$f'(x_1) \geq f'(x_2)$.

Se $f$ for duas vezes diferenciável em um intervalo, então $f$ é côncava se, e somente se,
$f''(x) \leq 0$.

Se, além disso,
$f''(x) < 0$,
então $f$ é estritamente côncava.

Assim, para funções suaves, a concavidade está associada a uma inclinação que diminui à medida que a variável aumenta.

Em macroeconomia, a concavidade é usada para representar \textbf{produto marginal decrescente}. 

No modelo de Solow-Swan, se a função de produção por trabalhador é $y=f(k)$, assume-se tipicamente que
$f'(k) > 0 \quad \text{e} \quad f''(k) < 0$.

A primeira condição diz que mais capital por trabalhador aumenta o produto por trabalhador; a segunda diz que esse aumento ocorre a taxas decrescentes. Logo, o produto marginal do capital é positivo, mas decrescente.

% ----------------------------------------------------------------------------
\newpage
\subsection{Condições de Inada}

\textbf{Condições de Inada} $\Rightarrow$ produtividade marginal explosiva quando o fator é escasso e nula quando o fator é abundante.

\medskip
\noindent\rule{\textwidth}{0.4pt}
\medskip

As \textbf{condições de Inada} são restrições sobre o comportamento assintótico das \textbf{produtividades marginais} de uma função de produção. Ou seja, descrevem o comportamento das derivadas da função de produção quando um fator tende a zero ou a infinito.

Considere uma função de produção $f(x)$ crescente, isto é, $f'(x)>0$, o que significa que o produto marginal é positivo; e côncava, isto é, $f''(x)<0$, o que significa que o produto marginal é decrescente. Nesse caso, as \textbf{condições de Inada} são dadas por:

$\lim_{x \to 0} f'(x)=\infty$
\quad e \quad
$\lim_{x \to \infty} f'(x)=0$.

Em palavras:

\begin{enumerate}
    \item quando o insumo se aproxima de zero, sua produtividade marginal tende ao infinito;
    \item quando o insumo cresce sem limite, sua produtividade marginal tende a zero.
\end{enumerate}

\medskip
\noindent\rule{\textwidth}{0.4pt}
\medskip

Para uma função de produção agregada $F(K,L)$, as condições de Inada podem ser formuladas em termos das produtividades marginais de cada fator.

Por exemplo, mantendo o trabalho $L$ fixo e variando o capital $K$, tem-se:

$\lim_{K \to 0} F_K(K,L)=\infty$
\quad e \quad
$\lim_{K \to \infty} F_K(K,L)=0$.

De forma análoga, mantendo o capital $K$ fixo e variando o trabalho $L$, tem-se:

$\lim_{L \to 0} F_L(K,L)=\infty$
\quad e \quad
$\lim_{L \to \infty} F_L(K,L)=0$.

\medskip
\noindent\rule{\textwidth}{0.4pt}
\medskip

Para a função Cobb-Douglas com uma variável:

$f(k)=k^\alpha$, \quad com $0<\alpha<1$.

Nesse caso,

$f'(k)=\alpha k^{\alpha-1}$,

de modo que

$\lim_{k \to 0} f'(k)=\infty$
\quad e \quad
$\lim_{k \to \infty} f'(k)=0$.

Logo, a Cobb-Douglas satisfaz as condições de Inada.

% ----------------------------------------------------------------------------
\newpage

\subsection{Função homogênea de grau $n$}

\textbf{Homogeneidade de grau $n$:} se todos os argumentos de uma função forem multiplicados por um mesmo escalar $\lambda$, o valor da função é multiplicado por $\lambda^n$.

$n = 1 \Rightarrow$ \textbf{retornos constantes de escala} $\Rightarrow$ $F(\lambda K,\lambda L)=\lambda F(K,L)$.

\medskip
\noindent\rule{\textwidth}{0.4pt}
\medskip

Uma função $F(K,L)$ é \textbf{homogênea de grau $n$} se, para todo $\lambda > 0$,
$F(\lambda K,\lambda L)=\lambda^n F(K,L)$.

A homogeneidade descreve como a função responde a uma variação proporcional de todos os seus argumentos.
Em particular:

\begin{enumerate}
    \item se $n=1$, a função é homogênea de grau 1 e apresenta \textbf{retornos constantes de escala};
    \item se $n>1$, a função apresenta \textbf{retornos crescentes de escala};
    \item se $n<1$, a função apresenta \textbf{retornos decrescentes de escala}.
\end{enumerate}

Se a função de produção agregada $F(K,L)$ é homogênea de grau 1, então

$F(\lambda K,\lambda L)=\lambda F(K,L)$,

isto é, ao multiplicar capital e trabalho por um mesmo fator, o produto também é multiplicado por esse fator.
Essa hipótese permite escrever a função de produção na forma intensiva (em termos por trabalhador, ou seja, dividindo por L). 

Se
$Y=F(K,L)$
e $F$ é homogênea de grau 1, então:

$\frac{Y}{L}=F\left(\frac{K}{L},1\right)$.

Definindo
$y=\frac{Y}{L}$ \quad e \quad $k=\frac{K}{L}$,
obtemos:

$y=f(k)$, $\quad \text{onde} \quad f(k)=F(k,1)$.

Essa transformação é central no modelo de Solow-Swan, pois permite estudar a dinâmica da economia em termos de produto por trabalhador e capital por trabalhador.
% ----------------------------------------------------------------------------
\newpage
\subsection{Teorema de Euler}

\textbf{Teorema de Euler:} se uma função é diferenciável e homogênea de grau $n$, então a soma das derivadas parciais ponderadas pelos respectivos argumentos é igual a $n$ vezes o valor da função.

\noindent\rule{\textwidth}{0.4pt}

Seja $f(x_1,\dots,x_m)$ uma função diferenciável e homogênea de grau $n$. Então,

$\sum_{i=1}^{m} \frac{\partial f}{\partial x_i}(x_1,\dots,x_m)\,x_i = n f(x_1,\dots,x_m)$.

\subsubsection{Aplicação ao modelo de Solow-Swan}

No modelo de Solow-Swan, a função de produção agregada é escrita como

$Y = F(K,L)$,

onde $K$ é o capital e $L$ é o trabalho.

Se $F(K,L)$ é diferenciável e homogênea de grau 1, isto é, apresenta retornos constantes de escala, então, pelo Teorema de Euler,

$F_K(K,L)\,K + F_L(K,L)\,L = F(K,L)$.

Como $Y = F(K,L)$, segue que

$F_K(K,L)\,K + F_L(K,L)\,L = Y$.

Aqui, $F_K(K,L)$ representa o produto marginal do capital e $F_L(K,L)$ representa o produto marginal do trabalho.

Sob concorrência perfeita, cada fator é remunerado de acordo com sua produtividade marginal. Assim, a remuneração por unidade de capital é

$r = F_K(K,L)$,

e a remuneração por unidade de trabalho é

$w = F_L(K,L)$.

Como a economia emprega $K$ unidades de capital e $L$ unidades de trabalho, a renda total do capital é

$rK = F_K(K,L)\,K$,

e a renda total do trabalho é

$wL = F_L(K,L)\,L$.

Substituindo essas expressões na identidade dada pelo Teorema de Euler, obtemos

$Y = rK + wL$.

Portanto, no modelo de Solow-Swan, sob retornos constantes de escala e concorrência perfeita, o produto total é exatamente exaurido pela remuneração dos fatores de produção. 

Assim, o Teorema de Euler fornece a base para a decomposição da renda entre capital e trabalho.


% ----------------------------------------------------------------------------
% Referências
% ----------------------------------------------------------------------------

%\newpage
%\renewcommand{\refname}{REFERÊNCIAS}
%\addcontentsline{toc}{section}{REFERÊNCIAS}
%\printbibliography

% ----------------------------------------------------------------------------
% Apêndice
% ----------------------------------------------------------------------------

%\newpage
%\appendix
%\section*{APÊNDICES} % título de abertura
%\addcontentsline{toc}{section}{APÊNDICES} % entra no sumário

% ------------------------------
% APÊNCICE A
% ------------------------------

%\section{Referência dos dados utilizados}
%\subsection{Produto}

% ------------------------------
% Final do documento
% ------------------------------

\end{document}
